% Minimal preamble for Open Logic Project sources (pinned commit 1e960beff9ed7835bf3e3f1335e21af3439cd107).
%
% OLP's own master preamble cannot run on the project host: TeX Live 2025 there lacks
% prooftree.sty, ntheorem.sty, extarrows.sty, mathdots.sty, isomath.sty and subfiles.sty.
% Following the project's own rule ("if the author preamble needs packages the toolchain lacks,
% do not invent macros -- use a minimal working preamble"), this file declares, verbatim in
% spirit and numbering behaviour, exactly the theorem-like environments OLP defines in
% sty/open-logic-book-envs.sty (thm/lem/prop/cor/defn/ex/prob/rem/note/case/axiom/conv),
% plus the shorthand accents OLP uses (`!!a`, `!!b`, `!!{...}`).
%
% Author text is untouched; only the packaging is ours.
\documentclass[10pt,a4paper]{article}

\usepackage{amsmath}
\usepackage{amssymb}
\usepackage{amsthm}
\usepackage{graphicx}
\usepackage{hyperref}
\usepackage{enumitem}
\usepackage[framemethod=default]{mdframed}

% ---- OLP shorthand accents (from sty/open-logic.sty) ----
% NB: \! is taken by LaTeX (negative thin space), so we only define the !!a / !!b families
% and the \olup helper used by them.
\newcommand*{\olup}[1]{\MakeUppercase{#1}}
\def\!!!a#1{\olup{#1}}
\def\!!a#1{\olup{#1}}
\def\!!b#1{\olup{#1}}

% ---- theorem-like environments, same family numbering as OLP ----
\theoremstyle{plain}
\newtheorem{thm}{Theorem}[section]
\newtheorem{lem}[thm]{Lemma}
\newtheorem{prop}[thm]{Proposition}
\newtheorem{cor}[thm]{Corollary}
\newtheorem{axiom}{Axiom}

\theoremstyle{definition}
\newtheorem{defn}[thm]{Definition}
\newtheorem{ex}[thm]{Example}
\newtheorem{prob}{Problem}[section]
\newtheorem{probd}{Problem}
\newtheorem{conv}[thm]{Convention}

\theoremstyle{remark}
\newtheorem{rem}[thm]{Remark}
\newtheorem{note}[thm]{Note}
\newtheorem{case}[thm]{Case}

% ---- environments OLP defines as wrappers ----
\newenvironment{probtag}[1]{\begin{prob}}{\end{prob}}
\newenvironment{explain}{}{}
\newenvironment{explain*}{}{}
\newenvironment{editorial}{}{}
\newenvironment{digress}{}{}
\newenvironment{digress*}{}{}
\newenvironment{history}{}{}
\newenvironment{reference}{}{}
\newcommand*{\olfileid}[3]{}
\newcommand*{\olsection}[1]{\section*{#1}}
\newcommand*{\olchapter}[3]{}
\newcommand*{\olpart}[2]{}
\newcommand*{\olimport}[1]{}
\newcommand*{\olinclude}[1]{}
\newcommand*{\olpath}[1]{}
\newcommand*{\ollangid}[1]{}
\newcommand*{\problemtag}[1]{}
\newcommand*{\problemauthor}[1]{}
\newcommand*{\ollabel}[1]{\label{#1}}
% \olref takes 0-3 OPTIONAL arguments plus one mandatory label:
%     \olref{lab}  |  \olref[from]{lab}  |  \olref[from][to]{lab}  |  \olref[a][b][c]{lab}
% A single-argument definition silently swallows "[sth]" and emits a broken reference
% (measured on the host: the text came out as "??iso]x"). The chain below consumes the
% brackets and cross-references only the mandatory argument.
\makeatletter
\newcommand{\olref}{\@ifnextchar[{\olref@i}{\olref@final}}
\def\olref@i[#1]{\@ifnextchar[{\olref@ii}{\olref@final}}
\def\olref@ii[#1]{\@ifnextchar[{\olref@iii}{\olref@final}}
\def\olref@iii[#1]{\olref@final}
\def\olref@final#1{\ref{#1}}
\makeatother
\newcommand*{\olnewpage}{\clearpage}

% ===================== OLP compatibility layer (1) =====================

\makeatletter

% ---------------- proof trees (OLP's prooftree is a bussproofs derivative) ----------------
\providecommand{\AxiomC}[1]{}
\providecommand{\UnaryInfC}[1]{}
\providecommand{\BinaryInfC}[1]{}
\providecommand{\TrinaryInfC}[1]{}
\providecommand{\QuaternaryInfC}[1]{}
\providecommand{\QuinaryInfC}[1]{}
\providecommand{\RightLabel}[1]{}
\providecommand{\LeftLabel}[1]{}
\providecommand{\doubleLine}{}
\providecommand{\noLine}{}
\providecommand{\DisplayProof}{}
\providecommand{\fCenter}{}
\newenvironment{prooftree}{}{}
\newenvironment{ltprooftree}{}{}
\newenvironment{oltableau}{}{}

% ---------------- logical notation ----------------
\providecommand{\lif}{\mathbin{\rightarrow}}
\providecommand{\lforall}{\mathord{\forall}}
\providecommand{\lexists}{\mathord{\exists}}
\providecommand{\ltrue}{\top}
\providecommand{\lfalse}{\bot}
\providecommand{\Ltrue}{\top}
\providecommand{\Lfalse}{\bot}
\providecommand{\Sequent}{\mathrel{\Rightarrow}}
\providecommand{\Proves}{\mathrel{\vdash}}
\providecommand{\Entails}{\mathrel{\models}}
\providecommand{\Subst}[3]{#1[#2/#3]}
\providecommand{\sFmla}[2]{#1(#2)}
\providecommand{\pFmla}[2]{#1(#2)}
\providecommand{\formula}{\relax}
\providecommand{\True}{\top}
\providecommand{\False}{\bot}
\providecommand{\Args}[1]{#1}
\providecommand{\Ts}{\ensuremath{\top}}
\providecommand{\Fi}{\ensuremath{\bot}}
\providecommand{\lfalsum}{\bot}

% ---------------- objects, structures, metatheory ----------------
\providecommand{\Obj}{}
\providecommand{\Struct}{\ensuremath{\mathfrak{A}}}
\providecommand{\Domain}{\ensuremath{|\mathfrak{A}|}}
\providecommand{\Assign}{\ensuremath{s}}
\providecommand{\Sat}[2]{\ensuremath{#1 \models #2}}
\providecommand{\mSat}[3]{\ensuremath{#1 \models #2[#3]}}
\providecommand{\mModel}[2]{\ensuremath{#1 \models #2}}
\providecommand{\Lang}{\ensuremath{\mathcal{L}}}
\providecommand{\Th}{\ensuremath{\mathrm{Th}}}
\providecommand{\Log}{\ensuremath{\mathrm{L}}}
\providecommand{\tuple}[1]{\langle #1 \rangle}
\providecommand{\num}[1]{\underline{#1}}
\providecommand{\cfind}[1]{\varphi_{#1}}
\providecommand{\Nat}{\ensuremath{\mathbb{N}}}
\providecommand{\Mult}{\ensuremath{\mathrm{Mult}}}
\providecommand{\Add}{\ensuremath{\mathrm{Add}}}
\providecommand{\Exp}{\ensuremath{\mathrm{Exp}}}
\providecommand{\Zero}{\ensuremath{0}}
\providecommand{\Succ}{\ensuremath{\mathrm{S}}}
\providecommand{\lambd}{\ensuremath{\lambda}}
\providecommand{\bredpar}{\mathrel{\rightarrow_\beta}}
\providecommand{\bcd}[1]{\ensuremath{\ulcorner #1 \urcorner}}
\providecommand{\Set}[1]{\ensuremath{\{#1\}}}
\providecommand{\indcase}[1]{\par\noindent\emph{Case #1}.\quad}
\providecommand{\indfrm}[1]{\par\noindent\emph{Form #1}.\quad}
\providecommand{\iftag}[1]{#1}
\providecommand{\tagitem}[1]{\item[#1]}
\providecommand{\Elim}[1]{}
\providecommand{\Intro}[1]{}
\providecommand{\LeftR}[1]{}
\providecommand{\RightR}[1]{}
\providecommand{\TAss}[1]{}
\providecommand{\TRule}[1]{}

\makeatother
\providecommand{\olpath}{.}

% ===================== OLP compatibility layer (2) =====================
% --- second compatibility block: notation and shorthands used by the frozen OLP units ---
\providecommand{\CutCS}{\mathrm{Cut}}
\providecommand{\LeftSubproofLabel}{\relax}
\providecommand{\RightSubproofLabel}{\relax}
\providecommand{\maeh}{\relax}
\providecommand{\cheight}{\relax}
\providecommand{\ande}{\mathrm{and}_E}
\providecommand{\andi}{\mathrm{and}_I}
\providecommand{\ore}{\mathrm{or}_E}
\providecommand{\ori}{\mathrm{or}_I}
\providecommand{\typeof}{\mathord{:}}
\providecommand{\pValue}[1]{#1}
\providecommand{\Sent}{\mathrm{Sent}}
\providecommand{\cardplus}{\mathbin{+}}
\providecommand{\cardtimes}{\mathbin{\cdot}}
\providecommand{\cardexpo}{\mathbin{\hat{\ }}}
\providecommand{\cardnless}{\not<}
\providecommand{\cardle}{\leq}
\providecommand{\funfromto}[3]{#1\colon #2 \to #3}
\providecommand{\funimage}[1]{\operatorname{im}(#1)}
\providecommand{\fdefined}{\mathord{\downarrow}}
\providecommand{\canonord}{\operatorname{ord}}
\providecommand{\ordtimes}{\mathbin{\cdot}}
\providecommand{\ordeq}[2]{#1 = #2}
\providecommand{\ordneq}[2]{#1 \neq #2}
\providecommand{\ordtype}[1]{\operatorname{ord}(#1)}
\providecommand{\ordsucc}{\mathord{+1}}
\providecommand{\ordunion}{\cup}
\providecommand{\onesphere}{S^1}
\providecommand{\rotationsgroup}{\mathrm{SO}(2)}
\providecommand{\equivrep}[1]{[#1]}
\providecommand{\rlexless}{\prec}
\providecommand{\ZFC}{\mathrm{ZFC}}
\providecommand{\ZF}{\mathrm{ZF}}
\providecommand{\Rat}{\mathbb{Q}}
\providecommand{\Real}{\mathbb{R}}
\providecommand{\PosInt}{\mathbb{Z}^{+}}
\providecommand{\Id}{\mathrm{id}}
\providecommand{\Bin}{\mathrm{Bin}}
\providecommand{\pto}{\rightharpoonup}
\providecommand{\dom}{\operatorname{dom}}
\providecommand{\ran}{\operatorname{ran}}
\providecommand{\len}{\operatorname{len}}
\providecommand{\concat}{\mathbin{{}^\frown}}
\providecommand{\card}[1]{\lvert #1 \rvert}
\providecommand{\Pow}{\mathcal{P}}
\providecommand{\Setabs}[2]{\{\, #1 \mid #2 \,\}}
\providecommand{\comp}{\circ}
\providecommand{\Char}{\chi}
\providecommand{\Prf}{\operatorname{Prf}}
\providecommand{\Prov}{\operatorname{Prov}}
\providecommand{\Proj}{\operatorname{proj}}
\providecommand{\Gn}{\operatorname{gn}}
\providecommand{\gn}{\operatorname{gn}}
\providecommand{\Value}[1]{\operatorname{val}(#1)}
\providecommand{\Deduce}{\mathrm{D}}
\providecommand{\Dual}{\mathrm{d}}
\providecommand{\Ax}{\mathrm{Ax}}
\providecommand{\pSat}[3]{#1 \models #2\,[#3]}
\providecommand{\pAssign}{\sigma}
\providecommand{\bforall}{\forall}
\providecommand{\bexists}{\exists}
\providecommand{\liff}{\mathbin{\leftrightarrow}}
\providecommand{\FV}{\operatorname{FV}}
\providecommand{\UnaryInf}{\ensuremath{\mid}}
\providecommand{\BinaryInf}{\ensuremath{\mid}}
\providecommand{\depth}{\operatorname{depth}}
\providecommand{\oliflabeldef}[3]{#1}
\providecommand{\TMblank}{\sqcup}
\providecommand{\TMstroke}{\mid}
\providecommand{\TMstay}{\mathrm{S}}
\providecommand{\TMright}{\mathrm{R}}
\providecommand{\TMleft}{\mathrm{L}}
\providecommand{\TMendtape}{\triangleright}
\providecommand{\umin}{\mu}
\providecommand{\Sum}{\sum}
\providecommand{\citep}[1]{\cite{#1}}
\providecommand{\citet}[1]{\cite{#1}}
\providecommand{\citeauthor}[1]{\cite{#1}}
\providecommand{\citeyear}[1]{\cite{#1}}

% ===================== OLP compatibility layer (3) =====================
% --- third compatibility block: remaining notation seen in the frozen OLP units ---
\providecommand{\iso}{\cong}
\providecommand{\nszero}{\mathrm{0}}
\providecommand{\nssucc}{\mathrm{S}}
\providecommand{\nsplus}{\mathbin{+}}
\providecommand{\nssless}{<}
\providecommand{\nsless}{<}
\providecommand{\PL}{\mathrm{PL}}
\providecommand{\Nec}{\mathord{\Box}}
\providecommand{\RK}{\mathrm{RK}}
\providecommand{\ST}{\mathrm{ST}}
\providecommand{\SSubst}[3]{#1[#2/#3]}
\providecommand{\subst}[3]{#1[#2/#3]}
\providecommand{\tagfalse}{\mathrm{F}}
\providecommand{\cutr}{\mathrm{Cut}}
\providecommand{\cardneq}{\neq}
\providecommand{\issucc}[1]{\operatorname{succ}(#1)}
\providecommand{\tup}[1]{\langle #1 \rangle}
\providecommand{\Var}{\mathrm{Var}}
\providecommand{\Bew}{\operatorname{Bew}}
\providecommand{\Con}{\operatorname{Con}}
\providecommand{\Cons}{\operatorname{Cons}}
\providecommand{\pDeduce}{\mathrm{D}}
\providecommand{\psat}{\models}
\providecommand{\den}[1]{\llbracket #1 \rrbracket}
\providecommand{\Sem}[2]{\llbracket #1 \rrbracket_{#2}}
\providecommand{\DomainOf}[1]{\lvert #1 \rvert}
\providecommand{\Form}{\mathrm{Form}}
\providecommand{\Terms}{\mathrm{Term}}
\providecommand{\Atom}{\mathrm{Atom}}
\providecommand{\FVars}[1]{\operatorname{FV}(#1)}
\providecommand{\fv}[1]{\operatorname{fv}(#1)}
\providecommand{\tuplep}[1]{\langle #1 \rangle}
\providecommand{\restrict}[2]{#1 \upharpoonright #2}
\providecommand{\sto}{\rightarrow}
\providecommand{\FmlaSet}[1]{\mathrm{Fmla}_{#1}}
\providecommand{\Thm}{\mathrm{Thm}}
\providecommand{\Eval}{\operatorname{val}}
\providecommand{\blank}{\sqcup}
\providecommand{\blankTape}{\sqcup}
\providecommand{\Consis}{\operatorname{Con}}
\providecommand{\drs}{\mathrm{d}}
\providecommand{\VRule}{\mid}
\providecommand{\Bet}{\mathrm{B}}
\providecommand{\Seg}[1]{\mathrm{Seg}_{#1}}
\providecommand{\Enums}{\mathrm{En}}
\providecommand{\Z}{\mathbb{Z}}
\providecommand{\N}{\mathbb{N}}
\providecommand{\abs}[1]{\lvert #1 \rvert}
\providecommand{\floor}[1]{\lfloor #1 \rfloor}
\providecommand{\ceil}[1]{\lceil #1 \rceil}

% ===================== OLP compatibility layer (4) =====================
% --- fourth compatibility block: harvested from the per-candidate compile logs ---
\providecommand{\beredpar}{\mathrel{\rightarrow_\beta}}
\makeatletter
\def\bered{\@ifnextchar[{\bered@i}{\mathrel{\rightarrow_\beta}}}
\def\bered@i[#1]{\mathrel{\rightarrow_\beta}}
\def\bred{\@ifnextchar[{\bred@i}{\mathrel{\rightarrow_\beta}}}
\def\bred@i[#1]{\mathrel{\rightarrow_\beta}}
\makeatother
\providecommand{\beredone}{\mathrel{=_\beta}}
\providecommand{\bredone}{\mathrel{=_\beta}}
\providecommand{\xredone}{\mathrel{=_\beta}}
\providecommand{\aconvone}{\mathrel{=_\alpha}}
\providecommand{\cardless}{<}
\providecommand{\tagrefs}[1]{#1}
\providecommand{\ordplus}{\mathbin{+}}
\providecommand{\funrestrictionto}[2]{#1 \upharpoonright #2}
\providecommand{\fundefined}{\mathord{\uparrow}}
\providecommand{\Trm}{\mathrm{Trm}}
\providecommand{\OProv}{\operatorname{Prov}}
\providecommand{\OPrf}{\operatorname{Prf}}
\providecommand{\aeq}{\mathrel{=}}
\providecommand{\mClass}[1]{[#1]}
\providecommand{\equivclass}[1]{[#1]}
\providecommand{\pheight}{\operatorname{ht}}
\providecommand{\phy}{\operatorname{ht}}
\providecommand{\Cut}{\mathrm{Cut}}
\providecommand{\QuantRank}{\operatorname{qr}}
\providecommand{\elemequiv}{\mathrel{\equiv}}
\providecommand{\PIso}{\mathrm{PIso}}
\providecommand{\Theory}{\mathrm{Th}}
\providecommand{\supstrict}{\subsetneq}
\providecommand{\trcl}{\operatorname{trcl}}
\providecommand{\cref}[1]{\ref{#1}}
\providecommand{\Discharge}{\mathrm{Dis}}
\providecommand{\Weakening}{\mathrm{Wk}}
\providecommand{\ZFminus}{\mathrm{ZF}^{-}}
\providecommand{\RK}{\mathrm{RK}}
\makeatletter
\renewcommand{\Dual}{\@ifnextchar[{\Dual@i}{\ifmmode\mathrm{d}\else$\mathrm{d}$\fi}}
\def\Dual@i[#1]{\ifmmode\mathrm{d}[#1]\else$\mathrm{d}[#1]$\fi}
\makeatother
\providecommand{\RightR}{\mathrm{R}}
\providecommand{\LeftR}{\mathrm{L}}
\providecommand{\Intro}{\mathrm{I}}
\providecommand{\Axiom}{\mathrm{Ax}}
\providecommand{\formula}[1]{#1}
\providecommand{\stared}{\ast}
\providecommand{\oset}[2]{\{\, #1 \mid #2 \,\}}
\providecommand{\setof}[2]{\{\, #1 \mid #2 \,\}}
\providecommand{\pow}[1]{\mathcal{P}(#1)}
\providecommand{\ordpair}[2]{\langle #1, #2 \rangle}
\providecommand{\ordtriple}[3]{\langle #1, #2, #3 \rangle}
\providecommand{\union}{\cup}
\providecommand{\inter}{\cap}
\providecommand{\seq}[1]{\langle #1 \rangle}
\providecommand{\SuchThat}{\mid}
\providecommand{\where}{\mid}
\providecommand{\TVar}{\mathrm{TVar}}
\providecommand{\Bool}{\mathrm{Bool}}
\providecommand{\Val}{\operatorname{Val}}
\providecommand{\FV}{\operatorname{FV}}
\providecommand{\TL}[1]{\mathrm{TL}}
\providecommand{\xstared}{\ast}

% ===================== OLP compatibility layer (5) =====================
% --- fifth compatibility block: rescues the quarantined OLP candidates ---
\providecommand{\ident}{\operatorname{id}}
\providecommand{\setrank}{\operatorname{rank}}
\providecommand{\cardfont}{\mathfrak{c}}
\providecommand{\Complement}{\overline}
\providecommand{\QR}{\mathrm{QR}}
\providecommand{\becd}{\ensuremath{\ulcorner}}
\providecommand{\disjointsum}{\mathbin{\oplus}}
\providecommand{\cardeq}[2]{#1 = #2}
\providecommand{\xred}{\mathrel{\rightarrow}}
\providecommand{\cardless}{<}
\providecommand{\cardplus}{\mathbin{+}}
\providecommand{\cardtimes}{\mathbin{\cdot}}
\providecommand{\cardexpo}{\mathbin{\hat{\ }}}
\providecommand{\cardneq}{\neq}
\providecommand{\cardle}{\leq}
\providecommand{\cardnless}{\not<}
% --- sixth compatibility block: names reported by real compile logs ---
\providecommand{\ORefut}{\operatorname{Refut}}
\providecommand{\ORProv}{\operatorname{Prov}}
\providecommand{\cutrank}{\operatorname{cr}}
\providecommand{\Refut}{\operatorname{Refut}}
\providecommand{\DischargeRule}{\mathrm{Dis}}
\providecommand{\Contraction}{\mathrm{C}}
\providecommand{\bottomAlignProof}{\relax}
\providecommand{\FalseCl}{\bot}
\providecommand{\red}{\mathrel{\rightarrow}}
\providecommand{\varAssign}{\sigma}
\providecommand{\defiff}{\mathrel{\leftrightarrow}}
\providecommand{\FalseInt}{\bot}
\providecommand{\fn}{\operatorname{fn}}
\providecommand{\stared}{\ast}
\providecommand{\dotted}{\cdot}
\providecommand{\Gamma}{\Gamma}
% --- sixth compatibility block: names reported by real compile logs ---
\providecommand{\ORefut}{\operatorname{Refut}}
\providecommand{\ORProv}{\operatorname{Prov}}
\providecommand{\cutrank}{\operatorname{cr}}
\providecommand{\Refut}{\operatorname{Refut}}
\providecommand{\DischargeRule}{\mathrm{Dis}}
\providecommand{\Contraction}{\mathrm{C}}
\providecommand{\bottomAlignProof}{\relax}
\providecommand{\FalseCl}{\bot}
\providecommand{\red}{\mathrel{\rightarrow}}
\providecommand{\varAssign}{\sigma}
\providecommand{\defiff}{\mathrel{\leftrightarrow}}
\providecommand{\FalseInt}{\bot}
\providecommand{\fn}{\operatorname{fn}}
\providecommand{\stared}{\ast}
\providecommand{\dotted}{\cdot}
\providecommand{\Gamma}{\Gamma}

% --- tag-macro semantic fix (2026-10-13, final): \iftag CHOOSES a branch and the OLP proof body is
% always the SECOND group (\iftag{tag}{body}, {tag}{body}{c}, {tag}{body}{c}{d}); a definition that
% prints the tag literal drops the proof body. Bodies contain blank lines, hence \long.
\makeatletter
\providecommand{\iftag}{}
\long\def\iftag#1#2{\@ifnextchar\bgroup{\iftag@three{#2}}{\iftag@stop{#2}}}
\long\def\iftag@three#1#2{\iftag@stop{#1}\@ifnextchar\bgroup{\iftag@eat}{\relax}}
\long\def\iftag@eat#1{}
\long\def\iftag@stop#1{#1}
\providecommand{\oliflabeldef}[2]{#1}
\long\def\oliflabeldef#1#2{\@ifnextchar\bgroup{\@gobble}{\iftag@stop{#2}}}
\providecommand{\tagrefs}[1]{}
\providecommand{\settag}[1]{}
\makeatother
\begin{document}

% ==== OLP file: content/incompleteness/incompleteness-provability/first-incompleteness-thm.tex ====
% Part: incompleteness
% Chapter: incompleteness-provability
% Section: first-incompleteness-thm



\olfileid{inc}{inp}{1in}

\olsection{The First Incompleteness Theorem}

We can now describe G\"odel's original proof of the first
incompleteness theorem. Let $\Th{T}$ be any computably axiomatized theory
in a language extending the language of arithmetic, such that $\Th{T}$
includes the axioms of $\Th{Q}$. This means that, in particular, $\Th{T}$
represents computable functions and relations.

We have argued that, given a reasonable coding of formulas and proofs
as numbers, the relation $\Prf[\Th{T}](x,y)$ is computable, where
$\Prf[\Th{T}](x,y)$ holds if and only if $x$ is the G\"odel number of
!!a{derivation} of the !!{formula} with G\"odel number~$y$
in~$\Th{T}$. In fact, for the particular theory that G\"odel had in
mind, G\"odel was able to show that this relation is primitive
recursive, using the list of 45 functions and relations in his
paper. The 45th relation, $x B y$, is just $\Prf[\Th{T}](x,y)$ for his
particular choice of~$\Th{T}$. Remember that where G\"odel uses the
word ``recursive'' in his paper, we would now use the phrase
``primitive recursive.''

Since $\Prf[\Th{T}](x,y)$ is computable, it is representable in $\Th{T}$. We
will use $\OPrf[\Th{T}](x,y)$ to refer to the formula that represents
it. Let $\OProv[\Th{T}](y)$ be the formula
$\lexists[x][\OPrf[\Th{T}](x,y)]$. This describes the 46th relation,
$\fn{Bew}(y)$, on G\"odel's list. As G\"odel notes, this is the only
relation that ``cannot be asserted to be recursive.''  What he
probably meant is this: from the definition, it is not clear that it
is computable; and later developments, in fact, show that it isn't.

Let $\Th{T}$ be an !!{axiomatizable} theory containing~$\Th{Q}$. Then
$\Prf[\Th{T}](x, y)$ is decidable, hence representable in~$\Th{Q}$ by
!!a{formula}~$\OPrf[\Th{T}](x, y)$. Let $\OProv[\Th{T}](y)$ be the formula we
described above. By the fixed-point lemma, there is a formula
$!G_\Th{T}$ such that $\Th{Q}$ (and hence $\Th{T}$) !!{derive}s
\begin{equation}
\ollabel{eqn:qpf}
!G_\Th{T} \liff \lnot \OProv[\Th{T}](\gn{!G_\Th{T}}).
\end{equation}
Note that $!G_\Th{T}$ says, in essence, ``$!G_\Th{T}$ is not
!!{derivable} in~$\Th{T}$.''

\begin{lem}\ollabel{lem:cons-G-unprov}
If $\Th{T}$ is a consistent, !!{axiomatizable} theory
extending~$\Th{Q}$, then $\Th{T} \Proves/ !G_\Th{T}$.
\end{lem}

\begin{proof}
Suppose $\Th{T}$ !!{derive}s $!G_\Th{T}$. Then there \emph{is}
!!a{derivation}, and so, for some number $m$, the relation $\Prf[\Th{T}](m,
\Gn{!G_\Th{T}})$ holds. But then $\Th{Q}$ !!{derive}s the sentence
$\OPrf[\Th{T}](\num m, \gn{!G_\Th{T}})$. So $\Th{Q}$ !!{derive}s
$\lexists[x][\OPrf[\Th{T}](x,\gn{!G_\Th{T}})]$, which is, by definition,
$\OProv[\Th{T}](\gn{!G_\Th{T}})$. By \olref{eqn:qpf}, $\Th{Q}$ !!{derive}s
$\lnot !G_\Th{T}$, and since $\Th{T}$ extends $\Th{Q}$, so
does~$\Th{T}$. We have shown that if $\Th{T}$ !!{derive}s $!G_\Th{T}$, then
it also !!{derive}s $\lnot !G_\Th{T}$, and hence it would be inconsistent.
\end{proof}

\begin{defn}
\ollabel{thm:oconsis-q}
A theory $\Th{T}$ is \emph{$\omega$-consistent} if the following holds: if
$\lexists[x][!A(x)]$ is any sentence and $\Th{T}$ !!{derive}s $\lnot
!A(\num 0)$, $\lnot !A(\num 1)$, $\lnot !A(\num 2)$, \dots then $\Th{T}$
does not prove $\lexists[x][!A(x)]$.
\end{defn}

Note that every $\omega$-consistent theory is also consistent. This
follows simply from the fact that if $\Th{T}$ is inconsistent, then
$\Th{T} \Proves !A$ for every~$!A$. In particular, if $\Th{T}$ is
inconsistent, it !!{derive}s both $\lnot !A(\num n)$ for every~$n$ and
also !!{derive}s~$\lexists[x][!A(x)]$. So, if $\Th{T}$ is
inconsistent, it is $\omega$-inconsistent. By contraposition, if
$\Th{T}$ is $\omega$-consistent, it must be consistent.

\begin{lem}\ollabel{lem:omega-cons-G-unref}
If $\Th{T}$ is an $\omega$-consistent, !!{axiomatizable} theory
extending~$\Th{Q}$, then $\Th{T} \Proves/ \lnot !G_\Th{T}$.
\end{lem}

\begin{proof}
We show that if $\Th{T}$ !!{derive}s $\lnot !G_\Th{T}$, then it is
$\omega$-inconsistent. Suppose $\Th{T}$ !!{derive}s $\lnot !G_\Th{T}$. If
$\Th{T}$ is inconsistent, it is $\omega$-inconsistent, and we are
done. Otherwise, $\Th{T}$ is consistent, so it does not !!{derive}
$!G_\Th{T}$ by \olref{lem:cons-G-unprov}. Since there is no
!!{derivation} of $!G_\Th{T}$ in $\Th{T}$, $\Th{Q}$ !!{derive}s
\[
\lnot \OPrf[\Th{T}](\num 0, \gn{!G_\Th{T}}), \lnot \OPrf[\Th{T}](\num 1,
\gn{!G_\Th{T}}), \lnot \OPrf[\Th{T}](\num 2, \gn{!G_\Th{T}}), \dots
\]
and so does~$\Th{T}$.  On the other hand, by \olref{eqn:qpf}, $\lnot
!G_\Th{T}$ is equivalent to
$\lexists[x][\OPrf[\Th{T}](x,\gn{!G_\Th{T}})]$. So $\Th{T}$ is
$\omega$-inconsistent.
\end{proof}

\begin{prob}
  Every $\omega$-consistent theory is consistent. Show that the
  converse does not hold, i.e., that there are consistent but
  $\omega$-inconsistent theories. Do this by showing that $\Th{Q} \cup
  \{\lnot !G_\Th{Q}\}$ is consistent but $\omega$-inconsistent.
\end{prob}

\begin{thm}
\ollabel{thm:first-incompleteness} Let $\Th{T}$ be any
$\omega$-consistent, !!{axiomatizable} theory extending~$\Th{Q}$. Then
$\Th{T}$ is not complete.
\end{thm}

\begin{proof}
  If $\Th{T}$ is $\omega$-consistent, it is consistent, so $\Th{T}
  \Proves/ !G_\Th{T}$ by \olref{lem:cons-G-unprov}.  By
  \olref{lem:omega-cons-G-unref}, $\Th{T} \Proves/ \lnot !G_\Th{T}$.
  This means that $\Th{T}$ is incomplete, since it !!{derive}s neither
  $!G_\Th{T}$ nor $\lnot !G_\Th{T}$.
\end{proof}


% ==== OLP file: content/incompleteness/incompleteness-provability/fixed-point-lemma.tex ====
% Part: incompleteness
% Chapter: incompleteness-provability
% Section: fixed-point-lemma



\olfileid{inc}{inp}{fix}

\olsection{The Fixed-Point Lemma}

\begin{explain}
The fixed-point lemma says that for any !!{formula}~$!B(x)$, there is
!!a{sentence}~$!A$ such that $\Th{T} \Proves !A \liff !B(\gn{!A})$,
provided $\Th{T}$ extends~$\Th{Q}$.  In the case of the liar sentence,
we'd want $!A$ to be equivalent (provably in~$\Th{T}$) to~``$\gn{!A}$
is false,'' i.e., the statement that $\Gn{!A}$ is the G\"odel number
of a false !!{sentence}. To understand the idea of the proof, it will
be useful to compare it with Quine's informal gloss of~$!A$ as,
``{}`yields a falsehood when preceded by its own quotation' yields a
falsehood when preceded by its own quotation.''  The operation of
taking an expression, and then forming a sentence by preceding this
expression by its own quotation may be called \emph{diagonalizing} the
expression, and the result its diagonalization. So, the
diagonalization of `yields a falsehood when preceded by its own
quotation' is ``{}`yields a falsehood when preceded by its own
quotation' yields a falsehood when preceded by its own quotation.''
Now note that Quine's liar sentence is not the diagonalization of
`yields a falsehood' but of `yields a falsehood when preceded by its
own quotation.' So the property being diagonalized to yield the liar
sentence itself involves diagonalization!{}

In the language of arithmetic, we form quotations of !!a{formula} with
one free variable by computing its G\"odel numbers and then
substituting the standard numeral for that G\"odel number into the
free variable. The diagonalization of~$!E(x)$ is $!E(\num{n})$, where
$n = \Gn{!E(x)}$. (From now on, let's abbreviate $\num{\Gn{!E(x)}}$ as
$\gn{!E(x)}$.)  So if $!B(x)$ is ``is a falsehood,'' then ``yields a
falsehood if preceded by its own quotation,'' would be ``yields a
falsehood when applied to the G\"odel number of its diagonalization.''
If we had a symbol~$\Obj{diag}$ for the function $\fn{diag}(n)$ which
computes the G\"odel number of the diagonalization of the !!{formula}
with G\"odel number~$n$, we could write $!E(x)$ as
$!B(\Obj{diag}(x))$. And Quine's version of the liar sentence would
then be the diagonalization of it, i.e., $!E(\gn{!E(x)})$ or
$!B(\Obj{diag}(\gn{!B(\Obj{diag}(x))}))$.  Of course, $!B(x)$ could
now be any other property, and the same construction would work. For
the incompleteness theorem, we'll take $!B(x)$ to be ``$x$~is not
!!{derivable} in~$\Th{T}$.'' Then $!E(x)$ would be ``yields
!!a{sentence} not !!{derivable} in~$\Th{T}$ when applied to the
G\"odel number of its diagonalization.''

To formalize this in~$\Th{T}$, we have to find a way to formalize
$\fn{diag}$. The function $\fn{diag}(n)$ is computable, in fact, it is
primitive recursive: if $n$ is the G\"odel number of a
formula~$!E(x)$, $\fn{diag}(n)$ returns the G\"odel number
of~$!E(\gn{!E(x)})$. (Recall, $\gn{!E(x)}$ is the standard numeral of
the G\"odel number of~$!E(x)$, i.e., $\num{\Gn{!E(x)}}$). If
$\Obj{diag}$ were a function symbol in $\Th{T}$ representing the
function $\fn{diag}$, we could take $!A$ to be the formula
$!B(\Obj{diag}(\gn{!B(\Obj{diag}(x))}))$. Notice that
\begin{align*}
\fn{diag}(\Gn{!B(\Obj{diag}(x))}) & = 
\Gn{!B(\Obj{diag}(\gn{!B(\Obj{diag}(x))}))} \\
& = \Gn{!A}.
\end{align*}
Assuming $\Th{T}$ can !!{derive}
\[
\Obj{diag}(\gn{!B(\Obj{diag}(x))}) = \gn{!A},
\]
it can !!{derive} $!B(\Obj{diag}(\gn{!B(\Obj{diag}(x))}))
\liff !B(\gn{!A})$. But the left hand side is, by
definition,~$!A$.

Of course, $\Obj{diag}$ will in general not be a function symbol of
$\Th{T}$, and certainly is not one of~$\Th{Q}$. But, since $\fn{diag}$
is computable, it is \emph{representable} in~$\Th{Q}$ by some formula
$!D_{\fn{diag}}(x,y)$. So instead of writing $!B(\Obj{diag}(x))$ we
can write $\lexists[y][(!D_{\fn{diag}}(x,y) \land !B(y))]$. Otherwise,
the proof sketched above goes through, and in fact, it goes through
already in~$\Th{Q}$.
\end{explain}

\begin{lem}
\ollabel{lem:fixed-point} Let $!B(x)$ be any formula with one free
variable~$x$. Then there is a sentence~$!A$ such that $\Th{Q} \Proves
!A \liff !B(\gn{!A})$.
\end{lem}


\begin{proof}
Given $!B(x)$, let $!E(x)$ be the formula
$\lexists[y][(!D_{\fn{diag}}(x,y) \land !B(y))]$ and let $!A$~be its
diagonalization, i.e., the formula $!E(\gn{!E(x)})$.

Since $!D_{\fn{diag}}$ represents $\fn{diag}$, and
$\fn{diag}(\Gn{!E(x)}) = \Gn{!A}$, $\Th{Q}$ can !!{derive}
\begin{align}
  & !D_{\fn{diag}}(\gn{!E(x)}, \gn{!A}) \ollabel{repdiag1} \\
  & \lforall[y][(!D_{\fn{diag}}(\gn{!E(x)},y) \lif
  \eq[y][\gn{!A}])]. \ollabel{repdiag2}
\end{align}
Now we show that $\Th{Q} \Proves !A \liff !B(\gn{!A})$. We argue
informally, using just logic and facts !!{derivable} in~$\Th{Q}$.

First, suppose~$!A$, i.e., $!E(\gn{!E(x)})$. Going back to the
definition of $!E(x)$, we see that $!E(\gn{!E(x)})$ just is
\[
\lexists[y][(!D_{\fn{diag}}(\gn{!E(x)},y) \land !B(y))].
\]
Consider such a~$y$. Since $!D_{\fn{diag}}(\gn{!E(x)},y)$, by
\olref{repdiag2}, $y = \gn{!A}$. So, from $!B(y)$ we
have~$!B(\gn{!A})$.

Now suppose $!B(\gn{!A})$. By \olref{repdiag1}, we have
\begin{align*}
& !D_{\fn{diag}}(\gn{!E(x)}, \gn{!A}) \land !B(\gn{!A}).
\intertext{It follows
that}
& \lexists[y][(!D_{\fn{diag}}(\gn{!E(x)},y) \land !B(y))].
\end{align*}
But that's just $!E(\gn{!E(x)})$, i.e.,~$!A$.
\end{proof}

\begin{digress}
You should compare this to the proof of the fixed-point lemma in
computability theory. The difference is that here we want to define a
\emph{statement} in terms of itself, whereas there we wanted to define
a \emph{function} in terms of itself; this difference aside, it is
really the same idea.
\end{digress}

\begin{prob} 
  !!^a{formula}~$!A(x)$ is a \emph{truth definition} if $\Th{Q}
  \Proves !B \liff !A(\gn{!B})$ for all !!{sentence}s~$!B$. Show that
  no !!{formula} is a truth definition by using the fixed-point lemma.
\end{prob}


% ==== OLP file: content/incompleteness/incompleteness-provability/godels-paper.tex ====
% Part: incompleteness
% Chapter: incompleteness-provability
% Section: godels-paper



\olfileid{inc}{inp}{gop}

\olsection{Comparison with G\"odel's Original Paper}

It is worthwhile to spend some time with G\"odel's 1931
paper. The introduction sketches the ideas we have just discussed.
Even if you just skim through the paper, it is easy to see what is
going on at each stage: first G\"odel describes the formal system $P$
(syntax, axioms, proof rules); then he defines the primitive recursive
functions and relations; then he shows that $x B y$ is primitive
recursive, and argues that the primitive recursive functions and
relations are represented in $\Th{P}$. He then goes on to prove the
incompleteness theorem, as above. In Section 3, he shows that one can
take the unprovable assertion to be a sentence in the language of
arithmetic. This is the origin of the $\beta$-lemma, which is what we
also used to handle sequences in showing that the recursive functions
are representable in $\Th{Q}$. G\"odel doesn't go so far to isolate a
minimal set of axioms that suffice, but we now know that $\Th{Q}$ will do
the trick.  Finally, in Section 4, he sketches a proof of the second
incompleteness theorem.


% ==== OLP file: content/incompleteness/incompleteness-provability/incompleteness-provability.tex ====
% Part: incompleteness
% Chapter: first-incompleteness



\olchapter{inc}{inp}{Incompleteness and Provability}

\olimport{introduction}

\olimport{fixed-point-lemma}

\olimport{first-incompleteness-thm}

\olimport{rosser-thm}

\olimport{godels-paper}

\olimport{provability-conditions}

\olimport{second-incompleteness-thm}

\olimport{lob-thm}

\olimport{tarski-thm}

\OLEndChapterHook


% ==== OLP file: content/incompleteness/incompleteness-provability/introduction.tex ====
% Part: incompleteness
% Chapter: incompleteness-provability
% Section: introduction



\olfileid{inc}{inp}{int}

\olsection{Introduction}

Hilbert thought that a system of axioms for a mathematical structure,
such as the natural numbers, is inadequate unless it allows one to
derive all true statements about the structure. Combined with his
later interest in formal systems of deduction, this suggests that he
thought that we should guarantee that, say, the formal systems we are
using to reason about the natural numbers is not only consistent, but
also \emph{complete}, i.e., every statement in its language is either
!!{derivable} or its negation is. G\"odel's first incompleteness theorem
shows that no such system of axioms exists: there is no complete,
consistent, !!{axiomatizable} formal system for arithmetic.  In fact,
no ``sufficiently strong,'' consistent, !!{axiomatizable} mathematical
theory is complete.

A more important goal of Hilbert's, the centerpiece of his program for
the justification of modern (``classical'') mathematics, was to find
finitary consistency proofs for formal systems representing classical
reasoning.  With regard to Hilbert's program, then, G\"odel's second
incompleteness theorem was a much bigger blow. The second
incompleteness theorem can be stated in vague terms, like the first
incompleteness theorem. Roughly speaking, it says that no sufficiently
strong theory of arithmetic can prove its own consistency. We will
have to take ``sufficiently strong'' to include a little bit more
than~$\Th{Q}$.

The idea behind G\"odel's original proof of the incompleteness theorem
can be found in the Epimenides paradox. Epimenides, a Cretan, asserted
that all Cretans are liars; a more direct form of the paradox is the
assertion ``this sentence is false.'' Essentially, by replacing truth
with !!{derivability}, G\"odel was able to formalize !!a{sentence}
which, in a roundabout way, asserts that it itself is not
!!{derivable}.  If that !!{sentence} were !!{derivable}, the theory
would then be inconsistent.  G\"odel showed that the negation of that
!!{sentence} is also not !!{derivable} from the system of axioms he was
considering. (For this second part, G\"odel had to assume that the
theory~$\Th{T}$ is what's called ``$\omega$-consistent.''
$\omega$-Consistency is related to consistency, but is a stronger
property.\footnote{That is, any $\omega$-consistent theory is
consistent, but not vice versa.} A few years after G\"odel, Rosser
showed that assuming simple consistency of~$\Th{T}$ is enough.)

The first challenge is to understand how one can construct !!a{sentence}
that refers to itself. For every !!{formula}~$!A$ in the language
of~$\Th{Q}$, let~$\gn{!A}$ denote the numeral corresponding
to~$\Gn{!A}$. Think about what this means: $!A$~is !!a{formula} in the
language of~$\Th{Q}$, $\Gn{!A}$~is a natural number, and $\gn{!A}$~is
a \emph{term} in the language of~$\Th{Q}$. So every !!{formula}~$!A$
in the language of~$\Th{Q}$ has a \emph{name},~$\gn{!A}$, which is a
term in the language of~$\Th{Q}$; this provides us with a conceptual
framework in which !!{formula}s in the language of~$\Th{Q}$ can ``say''
things about other !!{formula}s. The following lemma is known as the
fixed-point lemma.

\begin{lem}
Let $\Th{T}$ be any theory extending~$\Th{Q}$, and let $!B(x)$ be any
!!{formula} with only the variable~$x$ free. Then there is
!!a{sentence}~$!A$ such that $\Th{T} \Proves!A \liff !B(\gn{!A})$.
\end{lem}

The lemma asserts that given any property $!B(x)$, there is
!!a{sentence}~$!A$ that asserts ``$!B(x)$ is true of me,'' and
$\Th{T}$ ``knows'' this.

How can we construct such !!a{sentence}? Consider the following version
of the Epimenides paradox, due to Quine:
\begin{quote}
``Yields falsehood when preceded by its quotation'' yields falsehood
when preceded by its quotation.
\end{quote}
This sentence is not directly self-referential. It simply makes an
assertion about the syntactic objects between quotes, and, in doing
so, it is on par with sentences like
\begin{enumerate}
\item ``Robert'' is a nice name.
\item ``I ran.'' is a short sentence.
\item ``Has three words'' has three words.
\end{enumerate}
But what happens when one takes the phrase ``yields falsehood when
preceded by its quotation,'' and precedes it with a quoted version of
itself? Then one has the original sentence!{} In short, the sentence
asserts that it is false.


% ==== OLP file: content/incompleteness/incompleteness-provability/lob-thm.tex ====
% Part: incompleteness
% Chapter: incompleteness-provability
% Section: lob-thm



\olfileid{inc}{inp}{lob}

\olsection{L\"ob's Theorem}

The G\"odel sentence for a theory~$\Th{T}$ is a fixed point of $\lnot
\OProv[\Th{T}](y)$, i.e., !!a{sentence}~$!G$ such that
\[
\Th{T} \Proves \lnot \OProv[\Th{T}](\gn{!G}) \liff !G.
\]
It is not !!{derivable}, because if $\Th{T} \Proves !G$, (a) by !!{derivability}
condition~(1), $\Th{T} \Proves \OProv[\Th{T}](\gn{!G})$, and (b) $\Th{T}
\Proves !G$ together with $\Th{T} \Proves \lnot \OProv[\Th{T}](\gn{!G})
\liff !G$ gives $\Th{T} \Proves \lnot \OProv[\Th{T}](\gn{!G})$, and so
$\Th{T}$ would be inconsistent.  Now it is natural to ask about the
status of a fixed point of $\OProv[\Th{T}](y)$, i.e., !!a{sentence}~$!H$
such that
\[
\Th{T} \Proves \OProv[\Th{T}](\gn{!H}) \liff !H.
\]
If it were !!{derivable}, $\Th{T} \Proves \OProv[\Th{T}](\gn{!H})$ by
condition~(1), but the same conclusion follows if we apply modus
ponens to the equivalence above. Hence, we don't get that $\Th{T}$ is
inconsistent, at least not by the same argument as in the case of the
G\"odel sentence. This of course does not show that $\Th{T}$
\emph{does} !!{derive}~$!H$.

We can make headway on this question if we generalize it a bit. The
left-to-right direction of the fixed point equivalence,
$\OProv[\Th{T}](\gn{!H}) \lif !H$, is an instance of a general schema
called a \emph{reflection principle}: $\OProv[\Th{T}](\gn{!A}) \lif !A$.
It is called that because it expresses, in a sense, that $\Th{T}$ can
``reflect'' about what it can !!{derive}; basically it says, ``If $\Th{T}$
can !!{derive}~$!A$, then~$!A$ is true,'' for any~$!A$.  This is true for
sound theories only, of course, and this suggests that theories will
in general not !!{derive} every instance of it.  So which instances can a
theory (strong enough, and satisfying the !!{derivability} conditions)
!!{derive}?  Certainly all those where $!A$ itself is !!{derivable}. And that's
it, as the next result shows.

\begin{thm}\ollabel{thm:lob}
Let $\Th{T}$ be !!a{axiomatizable} theory extending $\Th{Q}$, and
suppose $\OProv[\Th{T}](y)$ is a formula satisfying conditions P1--P3 from
\olref[2in]{sec}. If $\Th{T}$ !!{derive}s $\OProv[\Th{T}](\gn{!A}) \lif !A$,
then in fact $\Th{T}$ !!{derive}s $!A$.
\end{thm}

Put differently, if $\Th{T} \Proves/ !A$, then $\Th{T} \Proves/
\OProv[\Th{T}](\gn{!A}) \lif !A$. This result is known as L\"ob's
theorem.

\begin{explain}
The heuristic for the proof of L\"ob's theorem is a clever proof that
Santa Claus exists. (If you don't like that conclusion, you are free
to substitute any other conclusion you would like.) Here it is:
\begin{enumerate}
\item Let $X$ be the sentence, ``If $X$ is true, then Santa Claus
  exists.''
\item Suppose $X$ is true.
\item Then what it says holds; i.e., we have: if $X$ is true, then
  Santa Claus exists.
\item Since we are assuming $X$ is true, we can conclude that
  Santa Claus exists, by modus ponens from (2) and~(3).
\item We have succeeded in deriving (4), ``Santa Claus exists,'' from
  the assumption~(2), ``$X$ is true.'' By conditional proof, we have
  shown: ``If $X$ is true, then Santa Claus exists.''
\item But this is just the sentence~$X$. So we have shown that $X$ is
  true.
\item But then, by the argument (2)--(4) above, Santa Claus exists.
\end{enumerate}
A formalization of this idea, replacing ``is true'' with ``is
!!{derivable},'' and ``Santa Claus exists'' with~$!A$, yields the proof of
L\"ob's theorem. The trick is to apply the fixed-point lemma to the
!!{formula}~$\OProv[\Th{T}](y) \lif !A$. The fixed point of that
corresponds to the sentence~$X$ in the preceding sketch.
\end{explain}

\begin{proof}[Proof of \olref{thm:lob}]
Suppose $!A$ is !!a{sentence} such that $\Th{T}$ !!{derive}s
$\OProv[\Th{T}](\gn{!A}) \lif !A$. Let $!B(y)$ be the !!{formula}~$\OProv[\Th{T}](y)
\lif !A$, and use the fixed-point lemma to find !!a{sentence}~$!D$
such that $\Th{T}$ !!{derive}s $!D \liff !B(\gn{!D})$. Then each of the
following is !!{derivable} in $\Th{T}$:
\begin{align}
  & !D \liff (\OProv[\Th{T}](\gn{!D}) \lif !A) \ollabel{L-1}\\
  & \qquad \text{$!D$ is a fixed point of~$!B(y)$}\notag \\
  & !D \lif (\OProv[\Th{T}](\gn{!D}) \lif !A) \ollabel{L-2}\\
  & \qquad\text{from \olref{L-1}}\notag\\
  & \OProv[\Th{T}](\gn{!D \lif (\OProv[\Th{T}](\gn{!D}) \lif !A)}) \ollabel{L-3}\\
  & \qquad \text{from \olref{L-2} by condition P1}\notag \\
  & \OProv[\Th{T}](\gn{!D}) \lif \OProv[\Th{T}](\gn{\OProv[\Th{T}](\gn{!D}) \lif !A})
  \ollabel{L-4}\\
  &\qquad \text{from \olref{L-3} using condition P2}\notag \\
  & \OProv[\Th{T}](\gn{!D}) \lif (\OProv[\Th{T}](\gn{\OProv[\Th{T}](\gn{!D})}) \lif \OProv[\Th{T}](\gn{!A})) \ollabel{L-5}\\
  &\qquad \text{from \olref{L-4} using P2 again} \notag\\
& \OProv[\Th{T}](\gn{!D}) \lif \OProv[\Th{T}](\gn{\OProv[\Th{T}](\gn{!D})}) \ollabel{L-6}\\
  & \qquad\text{by !!{derivability} condition P3} \notag\\
  & \OProv[\Th{T}](\gn{!D}) \lif \OProv[\Th{T}](\gn{!A}) \ollabel{L-7} \\
  &\qquad\text{from \olref{L-5} and \olref{L-6}}\notag\\
  & \OProv[\Th{T}](\gn{!A}) \lif !A \ollabel{L-8}\\
  &\qquad\text{by assumption of the theorem} \notag\\
  & \OProv[\Th{T}](\gn{!D}) \lif !A \ollabel{L-9}\\
  &\qquad\text{from \olref{L-7} and \olref{L-8}}\notag\\
  & (\OProv[\Th{T}](\gn{!D}) \lif !A) \lif !D \ollabel{L-10}\\
  & \qquad \text{from \olref{L-1}}\notag \\
  & !D \ollabel{L-11}\\
  & \qquad\text{from \olref{L-9} and \olref{L-10}}\notag \\
  & \OProv[\Th{T}](\gn{!D}) \ollabel{L-12}\\
  & \qquad\text{from \olref{L-11} by condition~P1}\notag \\
  & !A \qquad\qquad\text{from \olref{L-8} and \olref{L-12}}\notag
\end{align}
\end{proof}

With L\"ob's theorem in hand, there is a short proof of the second
incompleteness theorem (for theories having !!a{derivability} predicate
satisfying conditions P1--P3): if $\Th{T} \Proves
\OProv[\Th{T}](\gn{\lfalse}) \lif \lfalse$, then $\Th{T} \Proves \lfalse$.
If $\Th{T}$ is consistent, $\Th{T} \Proves/ \lfalse$. So, $\Th{T}
\Proves/ \OProv[\Th{T}](\gn{\lfalse}) \lif \lfalse$, i.e., $\Th{T} \Proves/
\OCon[\Th{T}]$.  We can also apply it to show that~$!H$, the fixed
point of $\OProv[\Th{T}](x)$, is !!{derivable}. For since
\begin{align*}
  \Th{T} & \Proves \OProv[\Th{T}](\gn{!H}) \liff !H\\
  \intertext{in particular}
    \Th{T} & \Proves \OProv[\Th{T}](\gn{!H}) \lif !H
\end{align*}
and so by L\"ob's theorem, $\Th{T} \Proves !H$.

% Going in the other direction, for homework I
% may ask you to work through a short proof of L\"ob's theorem, using
% the second incompleteness theorem instead of the fixed-point lemma.

\begin{prob}
Let $\Th{T}$ be a computably axiomatized theory, and
let $\OProv[\Th{T}]$ be !!a{derivability} predicate for $\Th{T}$. Consider the
following four statements:
\begin{enumerate}
\item If $T \Proves !A$, then $T \Proves \OProv[\Th{T}](\gn{!A})$.
\item $T \Proves !A \lif \OProv[\Th{T}](\gn{!A})$.
\item If $T \Proves \OProv[\Th{T}](\gn{!A})$, then $T \Proves !A$.
\item $T \Proves \OProv[\Th{T}](\gn{!A}) \lif !A$
\end{enumerate}
Under what conditions are each of these statements true?
\end{prob}


% ==== OLP file: content/incompleteness/incompleteness-provability/provability-conditions.tex ====
% Part: incompleteness
% Chapter: incompleteness-provability
% Section: provability-conditions



\olfileid{inc}{inp}{prc}

\olsection{The \usetoken{S}{derivability} Conditions for $\Th{PA}$}

Peano arithmetic, or $\Th{PA}$, is the theory extending $\Th{Q}$ with
induction axioms for all !!{formula}s. In other words, one adds to $\Th{Q}$
axioms of the form
\[
(!A(0) \land \lforall[x][(!A(x) \lif !A(x'))]) \lif \lforall[x][!A(x)]
\]
for every !!{formula}~$!A$. Notice that this is really a
\emph{schema}, which is to say, infinitely many axioms (and it turns
out that $\Th{PA}$ is {\em not} finitely axiomatizable). But since one
can effectively determine whether or not a string of symbols is an
instance of an induction axiom, the set of axioms for $\Th{PA}$ is
computable. $\Th{PA}$ is a much more robust theory than~$\Th{Q}$. For
example, one can easily prove that addition and multiplication are
commutative, using induction in the usual way. In fact, most finitary
number-theoretic and combinatorial arguments can be carried out
in~$\Th{PA}$.

Since $\Th{PA}$ is computably axiomatized, the !!{derivability}
predicate $\Prf[\Th{PA}](x,y)$ is computable and hence represented
in~$\Th{Q}$ (and so, in~$\Th{PA}$). As before, we will take
$\OPrf[\Th{PA}](x,y)$ to denote the formula representing the relation.
Let $\OProv[\Th{PA}](y)$ be the formula
$\lexists[x][\Prf[\Th{PA}](x,y)]$, which, intuitively says, ``$y$ is
!!{derivable} from the axioms of $\Th{PA}$.''  The reason we need a
little bit more than the axioms of $\Th{Q}$ is we need to know that
the theory we are using is strong enough to !!{derive} a few basic
facts about this !!{derivability} predicate. In fact, what we need are
the following facts:
\begin{enumerate}
\item[P1.] If $\Th{PA} \Proves !A$, then $\Th{PA} \Proves
  \OProv[\Th{PA}](\gn{!A})$.
\item[P2.] For all !!{formula}s $!A$ and $!B$,
  \[
  \Th{PA} \Proves \OProv[\Th{PA}](\gn{!A \lif !B}) \lif
  (\OProv[\Th{PA}](\gn{!A}) \lif \OProv[\Th{PA}](\gn{!B})).
  \]
\item[P3.] For every !!{formula}~$!A$,
  \[
  \Th{PA} \Proves \OProv[\Th{PA}](\gn{!A})
  \lif \OProv[\Th{PA}](\gn{\OProv[\Th{PA}](\gn{!A})}).
  \]
\end{enumerate}
The only way to verify that these three properties hold is to describe
the !!{formula} $\OProv[\Th{PA}](y)$ carefully and use the axioms of
$\Th{PA}$ to describe the relevant formal !!{derivation}s. Conditions (1)
and~(2) are easy; it is really condition~(3) that requires
work. (Think about what kind of work it entails \dots) Carrying out the
details would be tedious and uninteresting, so here we will ask you to
take it on faith that $\Th{PA}$ has the three properties listed
above. A reasonable choice of $\OProv[\Th{PA}](y)$ will also satisfy
\begin{enumerate}
\item[P4.] If $\Th{PA} \Proves \OProv[\Th{PA}](\gn{!A})$, then
  $\Th{PA} \Proves !A$.
\end{enumerate}
But we will not need this fact.

\begin{digress}
Incidentally, G\"odel was lazy in the same way we are
being now. At the end of the 1931 paper, he sketches the proof of the
second incompleteness theorem, and promises the details in a later
paper. He never got around to it; since everyone who understood the
argument believed that it could be carried out (he did not need to
fill in the details.)
\end{digress}


% ==== OLP file: content/incompleteness/incompleteness-provability/rosser-thm.tex ====
% Part: incompleteness
% Chapter: incompleteness-provability
% Section: rosser-thm



\olfileid{inc}{inp}{ros}

\olsection{Rosser's Theorem}

Can we modify G\"odel's proof to get a stronger result, replacing
``$\omega$-consistent'' with simply ``consistent''? The answer is
``yes,'' using a trick discovered by Rosser.  Rosser's trick is to use
a ``modified'' !!{derivability} predicate $\ORProv_T(y)$ instead of
$\OProv[\Th{T}](y)$.

\begin{thm}
\ollabel{thm:rosser}
Let $\Th{T}$ be any consistent, !!{axiomatizable} theory
extending $\Th{Q}$. Then $\Th{T}$ is not complete.
\end{thm}

\begin{proof}
Recall that $\OProv[\Th{T}](y)$ is defined as $\lexists[x][\OPrf[\Th{T}](x,
  y)]$, where $\OPrf[\Th{T}](x, y)$ represents the decidable relation which
holds iff $x$ is the G\"odel number of !!a{derivation} of the
!!{sentence} with G\"odel number~$y$. The relation that holds between
$x$ and~$y$ if $x$~is the G\"odel number of a \emph{refutation} of the
sentence with G\"odel number~$y$ is also decidable. Let $\fn{not}(x)$
be the primitive recursive function which does the following: if $x$
is the code of a formula $!A$, $\fn{not}(x)$ is a code of $\lnot
!A$. Then $\Refut[\Th{T}](x, y)$ holds iff $\Prf[\Th{T}](x, \fn{not}(y))$.  Let
$\ORefut[\Th{T}](x, y)$ represent it.  Then, if $\Th{T} \Proves \lnot !A$
and $\delta$ is a corresponding !!{derivation}, $\Th{Q} \Proves
\ORefut[\Th{T}](\gn{\delta}, \gn{!A})$.  We define $\ORProv[\Th{T}](y)$ as
\[
\lexists[x][(\OPrf[\Th{T}](x,y) \land \lforall[z][(z < x \lif \lnot
  \ORefut[\Th{T}](z,y))])].
\]
Roughly, $\ORProv[\Th{T}](y)$ says ``there is a proof of $y$ in $\Th{T}$,
and there is no shorter refutation of~$y$.''  Assuming $\Th{T}$ is
consistent, $\ORProv[\Th{T}](y)$ is true of the same numbers as
$\OProv[\Th{T}](y)$; but from the point of view of \emph{provability}
in~$\Th{T}$ (and we now know that there is a difference between truth
and provability!) the two have different properties. If $\Th{T}$ is
\emph{in}consistent, then the two do \emph{not} hold of the same
numbers!{} ($\ORProv[\Th{T}](y)$ is often read as ``$y$ is Rosser
provable.'' Since, as just discussed, Rosser provability is not some
special kind of provability---in inconsistent theories, there are
!!{sentence}s that are provable but not Rosser provable---this may be
confusing. To avoid the confusion, you could instead read it as
``$y$ is shmovable.'')

By the fixed-point lemma, there is a formula $!R_\Th{T}$ such that
\begin{equation}
  \Th{Q} \Proves !R_\Th{T} \liff \lnot \ORProv[\Th{T}](\gn{!R_\Th{T}}).
  \ollabel{RT}
\end{equation}
In contrast to the proof of \olref[1in]{thm:first-incompleteness},
here we claim that if $\Th{T}$ is consistent, $\Th{T}$ doesn't !!{derive}
$!R_\Th{T}$, and $\Th{T}$ also doesn't !!{derive} $\lnot !R_\Th{T}$. (In
other words, we don't need the assumption of $\omega$-consistency.)

First, let's show that $\Th{T} \Proves/ !R_{\Th{T}}$.  Suppose it did, so
there is !!a{derivation} of~$!R_{\Th{T}}$ from~$T$; let $n$ be its G\"odel
number. Then $\Th{Q} \Proves \OPrf[\Th{T}](\num{n}, \gn{!R_{\Th{T}}})$, since
$\OPrf[\Th{T}]$ represents $\Prf[\Th{T}]$ in~$\Th{Q}$. Also, for each $k < n$,
$k$ is not the G\"odel number of !!a{derivation} of $\lnot !R_{\Th{T}}$, since $\Th{T}$ is
consistent. So for each $k < n$, $\Th{Q} \Proves \lnot
\ORefut[\Th{T}](\num{k}, \gn{!R_{\Th{T}}})$. By \olref[req][min]{lem:less-nsucc},
$\Th{Q} \Proves \lforall[z][(z < \num{n} \lif \lnot \ORefut[\Th{T}](z,
  \gn{!R_{\Th{T}}}))]$. Thus,
\[
\Th{Q} \Proves \lexists[x][(\OPrf[\Th{T}](x,\gn{!R_{\Th{T}}}) \land \lforall[z][(z
    < x \lif \lnot \ORefut[\Th{T}](z,\gn{!R_{\Th{T}}}))])],
\]
but that's just $\ORProv[\Th{T}](\gn{!R_{\Th{T}}})$. By \olref{RT}, $\Th{Q}
\Proves \lnot !R_{\Th{T}}$. Since $\Th{T}$ extends $\Th{Q}$, also $\Th{T}
\Proves \lnot !R_{\Th{T}}$. We've assumed that $\Th{T} \Proves !R_{\Th{T}}$, so
$\Th{T}$ would be inconsistent, contrary to the assumption of the
theorem.

Now, let's show that $\Th{T} \Proves/ \lnot !R_{\Th{T}}$. Again, suppose it
did, and suppose $n$ is the G\"odel number of !!a{derivation}
of~$\lnot !R_{\Th{T}}$. Then $\Refut[\Th{T}](n, \Gn{!R_{\Th{T}}})$ holds, and since
$\ORefut[\Th{T}]$ represents $\Refut[\Th{T}]$ in $\Th{Q}$, $\Th{Q} \Proves
\ORefut[\Th{T}](\num{n}, \gn{!R_{\Th{T}}})$. We'll again show that $\Th{T}$ would
then be inconsistent because it would also !!{derive}~$!R_{\Th{T}}$.  Since
\begin{align*}
\Th{Q} & \Proves !R_{\Th{T}} \liff \lnot \ORProv[\Th{T}](\gn{!R_{\Th{T}}}), 
\intertext{and since $\Th{T}$ extends~$\Th{Q}$, it suffices to show that}
\Th{Q} & \Proves \lnot \ORProv[\Th{T}](\gn{!R_{\Th{T}}}).
\end{align*}
The !!{sentence} $\lnot
\ORProv[\Th{T}](\gn{!R_{\Th{T}}})$, i.e.,
\begin{align*}
  \lnot & \lexists[x][(\OPrf[\Th{T}](x,\gn{!R_{\Th{T}}}) \land
    \lforall[z][(z < x \lif \lnot \ORefut[\Th{T}](z,\gn{!R_{\Th{T}}}))])],
  \intertext{is logically equivalent to}
  & \lforall[x][(\OPrf[\Th{T}](x,\gn{!R_{\Th{T}}}) \lif
    \lexists[z][(z < x \land \ORefut[\Th{T}](z,\gn{!R_{\Th{T}}}))])].
\end{align*}
We argue informally using logic, making use of facts about what
$\Th{Q}$ !!{derive}s. Suppose $x$ is arbitrary and $\OPrf[\Th{T}](x,
\gn{!R_{\Th{T}}})$. We already know that $\Th{T} \Proves/ !R_{\Th{T}}$, and so for
every $k$, $\Th{Q} \Proves \lnot \OPrf[\Th{T}](\num{k}, \gn{!R_{\Th{T}}})$. Thus,
for every $k$ it follows that $\eq/[x][\num{k}]$. In particular, we
have (a) that $\eq/[x][\num{n}]$.  We also have $\lnot(\eq[x][\num{0}]
\lor \eq[x][\num{1}] \lor \dots \lor \eq[x][\num{n-1}])$ and so by
\olref[req][min]{lem:less-nsucc}, (b) $\lnot(x < \num{n})$. By
\olref[req][min]{lem:trichotomy}, $\num{n} < x$. Since $\Th{Q} \Proves
\ORefut[\Th{T}](\num{n}, \gn{!R_{\Th{T}}})$, we have $\num{n} < x \land
\ORefut[\Th{T}](\num{n}, \gn{!R_{\Th{T}}})$, and from that $\lexists[z][(z < x
\land \ORefut[\Th{T}](z,\gn{!R_{\Th{T}}}))]$. Since $x$ was arbitrary we get, as
required, that
\[
\lforall[x][(\OPrf[\Th{T}](x,\gn{!R_{\Th{T}}}) \lif
  \lexists[z][(z < x \land \ORefut[\Th{T}](z,\gn{!R_{\Th{T}}}))])].
\]
\end{proof}

\begin{prob}
Two sets $A$ and $B$ of natural numbers are said to be
\emph{computably inseparable} if there is no !!{decidable} set~$X$
such that $A \subseteq X$ and $B \subseteq \Complement{X}$
($\Complement{X}$ is the complement, $\Nat \setminus X$, of~$X$). Let
$\Th{T}$ be a consistent !!{axiomatizable} extension of $\Th{Q}$.
Suppose $A$ is the set of G\"odel numbers of !!{sentence}s provable
in~$\Th{T}$ and $B$ the set of G\"odel numbers of sentences
refutable in~$\Th{T}$. Prove that $A$ and~$B$ are computably
inseparable.
\end{prob}


% ==== OLP file: content/incompleteness/incompleteness-provability/second-incompleteness-thm.tex ====
% Part: incompleteness
% Chapter: incompleteness-provability
% Section: second-incompleteness-thm



\olfileid{inc}{inp}{2in}

\olsection{The Second Incompleteness Theorem}

How can we express the assertion that $\Th{PA}$ doesn't prove its own
consistency? Saying $\Th{PA}$ is inconsistent amounts to saying that
$\Th{PA} \Proves \eq[0][1]$. So we can take the consistency statement
$\OCon[\Th{PA}]$ to be the !!{sentence} $\lnot
\OProv[\Th{PA}](\gn{\eq[0][1]})$, and then the following theorem does
the job:

\begin{thm}
\ollabel{thm:second-incompleteness} 
Assuming $\Th{PA}$ is consistent, then $\Th{PA}$ does not !!{derive}
$\OCon[\Th{PA}]$.
\end{thm}

It is important to note that the theorem depends on the particular
representation of $\OCon[\Th{PA}]$ (i.e., the particular
representation of $\OProv[\Th{PA}](y)$). All we will use is that the
representation of $\OProv[\Th{PA}](y)$ satisfies the three
!!{derivability} conditions, so the theorem generalizes to any theory
with !!a{derivability} predicate having these properties.

It is informative to read G\"odel's sketch of an argument, since the
theorem follows like a good punch line. It goes like this. Let
$!G_\Th{PA}$ be the G\"odel sentence that we constructed in the proof
of \olref[1in]{thm:first-incompleteness}. We have shown ``If $\Th{PA}$
is consistent, then $\Th{PA}$ does not !!{derive} $!G_\Th{PA}$.'' If we
formalize this \emph{in} $\Th{PA}$, we have a proof of
\[
\OCon[\Th{PA}] \lif \lnot \OProv[\Th{PA}](\gn{!G_\Th{PA}}).
\]
Now suppose $\Th{PA}$ !!{derive}s $\OCon[\Th{PA}]$. Then it !!{derive}s $\lnot
\Prov[\Th{PA}](\gn{!G_\Th{PA}})$. But since $!G_\Th{PA}$ is a G\"odel
sentence, this is equivalent to $!G_\Th{PA}$. So $\Th{PA}$ !!{derive}s
$!G_\Th{PA}$.

But: we know that if $\Th{PA}$ is consistent, it doesn't !!{derive}
$!G_\Th{PA}$!{}  So if $\Th{PA}$ is consistent, it can't !!{derive}
$\OCon[\Th{PA}]$.

To make the argument more precise, we will let $!G_\Th{PA}$ be the
G\"odel sentence for~$\Th{PA}$ and use the !!{derivability} conditions
(P1)--(P3) to show that $\Th{PA}$ !!{derive}s $\OCon[\Th{PA}] \lif
!G_\Th{PA}$. This will show that $\Th{PA}$ doesn't !!{derive}
$\OCon[\Th{PA}]$. Here is a sketch of the proof, in~$\Th{PA}$. (For
simplicity, we drop the $\Th{PA}$ subscripts.)
\begin{align}
& !G \liff \lnot \OProv(\gn{!G}) \ollabel{G2-1}\\
& \qquad\text{$!G$ is a G\"odel sentence}\notag \\
& !G \lif \lnot \OProv(\gn{!G}) \ollabel{G2-2}\\
  & \qquad\text{from \olref{G2-1}} \notag\\
& !G \lif
  (\OProv(\gn{!G}) \lif \lfalse) \ollabel{G2-3}\\
  & \qquad\text{from \olref{G2-2} by logic}\notag\\
& \OProv(\gn{
    !G \lif
    (\OProv(\gn{!G}) \lif \lfalse)
  }) \ollabel{G2-4}\\
  & \qquad\text{by from \olref{G2-3} by condition P1} \notag\\
& \OProv(\gn{!G}) \lif
  \OProv(\gn{
    (\OProv(\gn{!G}) \lif \lfalse)
    }) \ollabel{G2-5}\\
  & \qquad\text{from \olref{G2-4} by condition P2} \notag\\
& \OProv(\gn{!G}) \lif (\OProv(\gn{\OProv(\gn{!G})}) \lif \OProv(\gn{\lfalse})) \ollabel{G2-6}\\
  & \qquad\text{from \olref{G2-5} by condition P2 and logic} \notag\\
& \OProv(\gn{!G}) \lif 
  \OProv(\gn{\OProv(\gn{!G})}) \ollabel{G2-7}\\
   & \qquad\text{by P3} \notag\\
& \OProv(\gn{!G}) \lif \OProv(\gn{\lfalse}) \ollabel{G2-8}\\
  & \qquad \text{from \olref{G2-6} and \olref{G2-7} by logic}\notag\\
& \OCon \lif \lnot \OProv(\gn{!G}) \ollabel{G2-9}\\
  & \qquad\text{contraposition of \olref{G2-8} and $\OCon \ident \lnot \OProv(\gn{\lfalse})$}\notag \\
& \OCon \lif !G \notag\\
  & \qquad\text{from \olref{G2-1} and \olref{G2-9} by logic}\notag
\end{align}
The use of logic in the above just elementary facts from propositional
logic, e.g., \olref{G2-3} uses $\Proves \lnot!A \liff (!A\lif
\lfalse)$ and \olref{G2-8} uses $!A \lif (!B \lif !C), !A \lif !B
\Proves !A \lif !C$. The use of condition~P2 in \olref{G2-5} and
\olref{G2-6} relies on instances of~P2, $\OProv(\gn{!A \lif !B}) \lif
(\OProv(\gn{!A}) \lif \OProv(\gn{!B}))$. In the first one, $!A \ident
!G$ and $!B \ident \OProv(\gn{!G}) \lif \lfalse$; in the second, $!A
\ident \OProv(\gn{G})$ and $!B \ident \lfalse$.

The more abstract version of the second incompleteness theorem is as follows:

\begin{thm}
\ollabel{thm:second-incompleteness-gen} Let $\Th{T}$ be any
consistent, !!{axiomatized} theory extending $\Th{Q}$ and let
$\OProv[\Th{T}](y)$ be any formula satisfying !!{derivability} conditions
P1--P3 for~$\Th{T}$. Then $\Th{T}$ does not !!{derive}~$\OCon[T]$.
\end{thm}

\begin{prob}
Show that $\Th{PA}$ !!{derive}s $!G_{\Th{PA}} \lif \OCon[\Th{PA}]$.
\end{prob}

\begin{digress}
The moral of the story is that no ``reasonable'' consistent theory for
mathematics can !!{derive} its own consistency statement. Suppose
$\Th{T}$ is a theory of mathematics that includes $\Th{Q}$ and
Hilbert's ``finitary'' reasoning (whatever that may be). Then, the
whole of $\Th{T}$ cannot !!{derive} the consistency statement of
$\Th{T}$, and so, a fortiori, the finitary fragment can't !!{derive}
the consistency statement of~$\Th{T}$ either. In that sense, there
cannot be a finitary consistency proof for ``all of mathematics.''

There is some leeway in interpreting the term ``finitary,'' and
G\"odel, in the 1931 paper, grants the possibility that something we
may consider ``finitary'' may lie outside the kinds of mathematics
Hilbert wanted to formalize. But G\"odel was being charitable; today,
it is hard to see how we might find something that can reasonably be
called finitary but is not formalizable in, say, $\Th{ZFC}$,
Zermelo--Fraenkel set theory with the axiom of choice.
\end{digress}


% ==== OLP file: content/incompleteness/incompleteness-provability/tarski-thm.tex ====
% Part: incompleteness
% Chapter: incompleteness-provability
% Section: lob-thm



\olfileid{inc}{inp}{tar}

\olsection{The Undefinability of Truth}

The notion of \emph{definability} depends on having a formal semantics
for the language of arithmetic.  We have described a set of formulas
and sentences in the language of arithmetic. The ``intended
interpretation'' is to read such sentences as making assertions about
the natural numbers, and such an assertion can be true or false. Let
$\Struct{N}$ be the !!{structure} with domain $\Nat$ and the standard
interpretation for the symbols in the language of arithmetic.  Then
$\Sat{N}{!A}$ means ``$!A$ is true in the standard interpretation.''

\begin{defn}
A relation $R(x_1,\dots,x_k)$ of natural numbers is \emph{definable}
in $\Struct{N}$ if and only if there is a formula $!A(x_1,\dots,x_k)$
in the language of arithmetic such that for every $n_1,\dots,n_k$,
$R(n_1,\dots,n_k)$ if and only if $\Sat{N}{!A(\num n_1,\dots,\num
  n_k)}$.
\end{defn}

Put differently, a relation is definable in $\Struct{N}$ if and
only if it is representable in the theory $\Th{TA}$, where $\Th{TA} =
\Setabs{!A}{\Sat{N}{!A}}$ is the set of true sentences of
arithmetic. (If this is not immediately clear to you, you should go
back and check the definitions and convince yourself that this is the
case.)

\begin{lem}
Every computable relation is definable in~$\Struct{N}$.
\end{lem}

\begin{proof}
It is easy to check that the formula representing a relation in
$\Th{Q}$ defines the same relation in $\Struct{N}$. 
\end{proof}

Now one can ask, is the converse also true?  That is, is every
relation definable in~$\Struct{N}$ computable? The answer is no. For
example:

\begin{lem}
The halting relation is definable in $\Struct{N}$.
\end{lem}

\begin{proof}
Recall that the Kleene normal form theorem states that every partial
computable function~$f$ has an index~$e$ such that $f(x) =
U(\umin{s}{T(e,x,s)})$ for all $x \in \Nat$, where $U$ and $T$ are
primitive recursive and therefore total. Thus, $f(x)$ is defined
(i.e., the computation halts) iff there is an~$s$ such that $T(e,x,s)$
holds.

Now let $H$ be the halting relation, i.e.,
\[
H = \Setabs{\tuple{e,x}}{\lexists[s][T(e, x, s)]}.
\]
Let $!D_T$ define $T$ in $\Struct{N}$. Then
\[
H = \Setabs{\tuple{e,x}}{\Sat{N}{\lexists[s][!D_T(\num e, \num x, s)]}},
\]
so $\lexists[s][!D_T(z, x, s)]$ defines~$H$ in $\Struct{N}$. 
\end{proof}

\begin{prob}
Show that $Q(n) \defiff n \in \Setabs{\Gn{!A}}{\Th{Q} \Proves !A}$ is
  definable in arithmetic.
\end{prob}

What about $\Th{TA}$ itself? Is it definable in arithmetic? That
is: is the set $\Setabs{\Gn{!A}}{\Sat{N}{!A}}$ definable in
arithmetic? Tarski's theorem answers this in the negative.

\begin{thm}
\ollabel{thm:tarski}
The set of true !!{sentence}s of arithmetic is not definable in arithmetic.
\end{thm}

\begin{proof} 
Suppose $!D(x)$ defined it, i.e., $\Sat{N}{!A}$ iff
$\Sat{N}{!D(\gn{!A})}$. By the fixed-point lemma, there is a formula
$!A$ such that $\Th{Q} \Proves !A \liff \lnot !D(\gn{!A})$, and hence
$\Sat{N}{!A \liff \lnot !D(\gn{!A})}$. But then $\Sat{N}{!A}$ if and
only if $\Sat{N}{\lnot !D(\gn{!A})}$, which contradicts the fact that
$!D(y)$ is supposed to define the set of true statements of
arithmetic.  
\end{proof}

Tarski applied this analysis to a more general philosophical notion of
truth. Given any language $L$, Tarski argued that an adequate notion
of truth for $L$ would have to satisfy, for each sentence $X$,
\begin{quote}
`$X$' is true if and only if $X$.
\end{quote}
Tarski's oft-quoted example, for English, is the sentence
\begin{quote}
`Snow is white' is true if and only if snow is white.
\end{quote}
However, for any language strong enough to represent the diagonal
function, and any linguistic predicate $T(x)$, we can construct a
sentence $X$ satisfying ``$X$ if and only if not $T(\text{`$X$'})$.''
Given that we do not want a truth predicate to declare some sentences
to be both true and false, Tarski concluded that one cannot specify a
truth predicate for all sentences in a language without, somehow,
stepping outside the bounds of the language. In other words, a the
truth predicate for a language cannot be defined in the language
itself.
